\begin{Problem}[Exercise 6.1: Spectrum-generating algebra]
Let us consider a three-site symmetric exclusion process $(w_L = w_R = 1)$
with closed boundaries as shown in the figure. Its Liouvillian is given by
\[
\mathcal{L} = \mathcal{L}_{12} + \mathcal{L}_{23}
\]
where
\[
X := \mathcal{L}_{12} = \mathcal{L}^{(2)} \otimes \mathbbm{1},
\qquad
Y := \mathcal{L}_{23} = \mathbbm{1} \otimes \mathcal{L}^{(2)}.
\]

$\mathcal{L}$ has the eigenvalues
\[
\{0,0,0,0,1,1,3,3\}
\]
(on quartett and two doublets). Our aim is to compute the eigenvalues of
$\mathcal{L}$ solely by algebraic methods without using any representation,
in a similar way as one solves the harmonic oscillator in terms of $a,a^\dagger$.

\begin{enumerate}
    \item[(a)] Verify that the matrices $X$ and $Y$ obey the so-called
    Temperley--Lieb algebra
    \[
    \text{Idempotence}
    \qquad
    X^2 = 2X,
    \qquad
    Y^2 = 2Y
    \]
    \[
    \text{Braid group property}
    \qquad
    XYX = X,
    \qquad
    YXY = Y.
    \]

    Please use from now on exclusively the symbols $X$ and $Y$ and the four
    algebraic relations given above. Do not use their explicit matrix
    representation.

    \item[(b)] Specify the monomials of the algebra, i.e.\ the elementary words
    formed by the letters $X$ and $Y$ that cannot be reduced by means of the
    algebraic relations.

    \item[(c)] Joining two monomials by concatenation and applying the algebraic
    relations one obtains another monomial. List all possible results in a
    table.

    \item[(d)] Now consider a polynomial, i.e.\ a linear combination of all
    monomials with arbitrary coefficients. Apply the Liouvillian to this
    polynomial and find the resulting polynomial, showing how the Liouvillian
    acts in the linear space of polynomials.

    \item[(e)] Determine the eigenpolynomials of
    \[
    \mathcal{L} = X + Y,
    \]
    i.e.\ those polynomials which are mapped by $\mathcal{L}$ onto themselves
    up to a factor. Calculate the corresponding eigenvalues and compare them
    with the spectrum given above.
\end{enumerate}
\end{Problem}
\begin{proof}
	\begin{parts}
	\item The symmetric exclusion process is defined in the lecture to act on two adjacent sites with the matrix representation
		\[
			\mathcal{L}^{(2)} = \begin{pmatrix} 0 & 0 & 0 & 0 \\ 0 & w_L & -w_R & 0 \\ 0 & -w_L & w_R & 0 \\ 0 & 0 & 0 & 0 \end{pmatrix} 
		\] 
		in the canonical basis $(\emptyset\emptyset, \emptyset A, A\emptyset, A A)$.

		To show idempotence of $X$ and $Y$, we only need to show that $(\mathcal{L}^{(2)})^2=2\mathcal{L}^{(2)}$, because then we would have
		\begin{align*}
			X^2&= \mathcal{L}^{(2)}\mathcal{L}^{(2)}\otimes \mathbbm{1}=2\mathcal{L}^{(2)}\otimes \mathbbm{1}=2X\\
			Y^2&= \mathbbm{1}\otimes\mathcal{L}^{(2)}\mathcal{L}^{(2)}=2\mathbbm{1}\otimes\mathcal{L}^{(2)}=2Y
		\end{align*}
		Additionally, since $\mathcal{L}^{(2)}$ has the block diagonal representation $\mathcal{L}^{(2)}=\text{diag}(0, L, 0)$ where
		\[
			L = \begin{pmatrix} w_L & -w_R \\ -w_L & w_R \end{pmatrix} 
		,\]
		we only need to show it for the $2\times 2$ matrix $L$, since clearly the $1\times 1$ matrix 0 is idempotent. Thus, we compute
		\[
		L^2 =\left(
\begin{array}{cc}
 w_L w_R+w_L^2 & -w_L w_R-w_R^2 \\
 -w_L w_R-w_L^2 & w_L w_R+w_R^2 \\
\end{array}
\right) 
		.\] 
		Substituing $w_L=w_R=1$, we find
		\[
			L^2 = 2 \begin{pmatrix} 1 & -1 \\ -1 & 1 \end{pmatrix}=2L 
		\] 
		as desired.
	\item Clearly, $X$ and $Y$ are themselves words. We consider now how we can add letters to these single letter words to extend them. We cannot add an $X$ to $X $, because that would then be reducible by idempotence, and similarly we cannot add $Y$ to $Y$. Thus, we append the other letter to them to find that $XY$ and $YX$ are monomials. 

		Now, again, we cannot add an $X$ to the left of $XY$, because it would be reducible by idempotence. We cannot add a $Y$ to its right, because it would then be reducible by the braid group property. The opposite argument holds with appending $Y$s to either side, and the same to $YX$. Thus, we cannot construct any three letter words that are irreducible.

		We conclude that the only monomials are $X,~Y,~XY,$ and $YX$.
	\item There are 4 monomials, resulting in $4\times 4 = 16$ possible products. We compute all possible results
		\begin{align*}
			XX &= X^2 = 2X \\
			XY &= XY \\
			X(XY) &= X^2 Y = 2XY \\
			X(YX) &= XYX = X \\
			YX &= YX \\
			YY &= Y^2 = 2Y \\
			Y(XY) &= YXY = Y \\
			YYX &= Y^2 X = 2YX \\
			(XY)X &= XYX = X \\
			(XY)Y &= XY^2 = 2XY \\
			(XY)(XY) &=  (XYX)Y = XY \\
			(XY)(YX) &= XY^2X=2XYX = 2X\\
			(YX)X &= YX^2 = 2YX \\
			(YX)Y &= YXY = Y \\
			(YX)(XY) &= YX^2Y = 2YXY = 2Y\\
			(YX)(YX) &= (YXY)X = YX  
		\end{align*}
		As requested, we summarise this in a table
\begin{center}
\renewcommand\arraystretch{1.3}
\setlength\doublerulesep{0pt}
\begin{tabular}{c|c|c|c|c|}
$\cdot$ & $X$ & $Y$ & $XY$ & $YX$ \\
\hline
	$X$ & $2X$ & $XY$ & $2XY$ & $X$ \\ 
\hline
$Y$ & $YX$ & $2Y$ & $Y$ & $2YX$ \\ 
\hline
$XY$ & $X$ & $2XY$ & $XY$ & $2X$ \\ 
\hline
$YX$ & $2YX$ & $Y$ & $2Y$ & $YX$ \\ 
\hline
\end{tabular}
\end{center}
The table is to be read as row $\times$ column.
\item We write a polynomial as, given 4 coefficients $c_1,c_2,c_3,c_4$ in the underlying field,
	\[
		P[X,Y](c_1,c_2,c_3,c_4) = c_1 X + c_2 Y + c_3 XY + c_4 YX
	.\] 
	We then apply the Liouvillian to this polynomial
	\begin{align*}
		&~\mathcal{L}(c_1 X + c_2 Y + c_3 XY + c_4 YX)\\
		=&~(X+Y)(c_1 X + c_2 Y + c_3 XY + c_4 YX)\\
		=&~c_1 X^2 + c_2 XY + c_3 XXY + c_4 XYX\\
		 &~+ c_1 YX + c_2 Y^2 + c_3 YXY + c_4 YYX\\
		=&~2c_1 X + c_2 XY + 2c_3 XY + c_4 X\\
		 &~+ c_1 YX + 2 c_2 Y + c_3 Y + 2c_4 YX\\
		=&~(2c_1+c_4)X + (2c_2+c_3)Y + (c_2+2c_3)XY + (c_1+2c_4)YX
	\end{align*}
	Thus, in the linear space of polynomials, which we identify with $(c_1,c_2,c_3,c_4)\cong \R^3$, the Liouvillian acts as
	\[
		\mathcal{L}:\R^4 \to \R^4,\qquad \begin{pmatrix} c_1 \\ c_2 \\ c_3 \\c_4 \end{pmatrix} \mapsto \begin{pmatrix} 2 & 0 & 0 & 1\\
	0 & 2 & 1 & 0 \\
0 & 1 & 2 & 0 \\
1 & 0 & 0 & 2\end{pmatrix} \begin{pmatrix} c_1 \\ c_2 \\ c_3 \\ c_4 \end{pmatrix} 
	.\] 
	\end{parts}
	We diagonalise this matrix, again, using python
	\begin{lstlisting}[language=Python]
		import numpy as np
		L = np.array([[2,0,0,1],[0,2,1,0],[0,1,2,0],[1,0,0,2]])
		np.linalg.eig(L)
	\end{lstlisting}
	which yields the eigenpolynomials and their associated eigenvalues $\lambda$:
	\begin{align*}
		\lambda=3:&~ X + YX\\
		\lambda=3:&~Y+XY\\
		\lambda=1:&~X - YX\\
		\lambda=1:&~-Y + XY
	\end{align*}
\end{proof}
